\section{Detection design and simulation evidence}\label{sec:sim}

No public dataset links Thai exchange accounts to criminal outcomes, and the data that do exist cannot be shared. To
learn how the estimators of \cref{sec:framework} behave before such data are available, I simulate exchange-years in
which the truth is known. The simulator is a method-validation device: it is built to make detection \emph{hard}, its
parameters are stylised rather than calibrated to any proprietary dataset, and none of its outputs is evidence about
real markets.

\subsection{The simulated economy}\label{sec:simdesign}

Each draw simulates 8,000 accounts on one licensed platform over 365 days. An account generates \emph{cycles}: a baht
deposit, the purchase it funds and, possibly, an on-chain withdrawal. Honest accounts belong to five populations that
embody the legitimate motives of \cref{sec:motives}; \cref{tab:populations} summarises them and \cref{app:simparams}
lists every parameter.

\begin{table}[!htbp]
\centering\small
\begin{threeparttable}
\caption{Populations in the simulated exchange-year}
\label{tab:populations}
\begin{tabularx}{\textwidth}{L{2.7cm}cccL{6.4cm}}
\toprule
Population & Share & \makecell{Median\\deposit (THB)} & \makecell{Median\\exit} & Distinguishing behaviour \\
\midrule
Domestic traders & 53\% & 4,000 & 3 days & Mostly non-stablecoin purchases; rarely withdraw \\
Offshore bridge traders & 21\% & 20,000 & 20 min & 90\% USDT; withdraw to exchange hot wallets; 15\% are arbitrageurs exiting in about 5 min \\
Cross-border remitters & 10\% & 6,000 & 6 hours & Small regular USDT purchases forwarded to P2P merchants abroad, 10\% of whom are tainted \\
Holders & 13\% & 15,000 & 30 days & Buy USDT and keep it on the platform \\
Honest P2P merchants & 3\% & 30,000 & 15 min & Many counterparties (70\% of deposits from new bank accounts); fast exits \\
\midrule
Mules & 0.25--4\% & structured & 3--90 min & Draw their own tradecraft; 40\% camouflaged on rented aged accounts with pre-layered funding \\
\bottomrule
\end{tabularx}
\begin{tablenotes}[flushleft]\footnotesize
\item \emph{Notes:} Shares of honest accounts are normalised to sum to one; mules are added at the stated share of all
accounts. 35\% of honest deposits intended for 50,000--62,500 baht relocate just below 50,000 baht in every population
(\cref{as:a1}). Amounts heap at multiples of 1,000 (30\% of deposits) and 5,000 baht (12\%).
\end{tablenotes}
\end{threeparttable}
\end{table}

Mules differ from honest accounts on several dimensions at once, but never on all of them. Each mule draws its own
exit speed (log-uniform between 3 and 90 minutes), its propensity to structure deposits into the 40,000--50,000-baht
band (20--90\%) and its share of deposits from fresh bank accounts (40--95\%). Forty per cent are \emph{camouflaged}:
they operate rented accounts opened up to three years earlier, receive pre-layered funds from one or two bank accounts,
structure little, exit no faster than 30 minutes and route half their withdrawals through exchange hot wallets.

The transaction graph contains 60 flagged addresses, 400 layering wallets, 20 exchange hot wallets and 200 peer-to-peer
merchants, plus a personal wallet per account and consolidation wallets for mules. Blacklists are incomplete by
design: half of the layering wallets pay a flagged address directly, a fifth do so through another layering wallet and
the rest belong to infrastructure nobody has flagged. Forty per cent of layering wallets cash out through exchange hot
wallets, which therefore sit two undirected hops from a flagged address---the premise of \cref{prop:direction}.

\subsection{Estimators and evaluation}

The bunching-difference test uses 1,000-baht bins between 10,000 and 110,000 baht, an exclusion window from 45,000 to
65,000 baht, a Poisson counterfactual that is cubic in $\log(z/T)$ with multiplicative round-number effects, and a
Poisson bootstrap with 200 replications. The dwell score is an account's mean posterior probability of belonging to the
fastest component of a BIC-selected Gaussian mixture on log dwell times. Exposure is \cref{eq:exposure} with $k=2$ on
the undirected graph and $k=3$ on the directed graph. The red-flag composite averages the percentile ranks of eight
indicators adapted from the FATF's red flags \citep{fatf2020redflags}: account youth, funding fan-in, stablecoin share,
fast exits, deposits near the threshold, unlabelled destinations, daytime activity and directed exposure. The
isolation forest uses the same indicators. ``Behaviour'' variants drop the exposure indicator, mimicking an exchange
without blockchain-analytics data.

Because mules are rare, I report the area under the ROC curve (AUC) and \emph{precision at $k$}, the share of mules
among the $k$ highest-scored accounts when $k$ equals the number of mules. Precision at $k$ is what a compliance team
with capacity to review $k$ accounts experiences; AUC averages over thresholds nobody uses. For tests, I report size
(the rejection rate when no mules exist) and power.

\subsection{Results}

\paragraph{Structuring is visible as a difference}
\Cref{fig:bunching} shows one simulated exchange-year with 2\% mules. Both channels bunch below the threshold because
every population avoids the scan: the responding share is 0.34 for deposits funding other assets, close to the
true honest response of 0.35, and 0.51 for deposits funding USDT. The difference, 0.17 (s.e.\ 0.06), rejects the null
of no structuring with a one-sided $p$-value of 0.001. Across the Monte Carlo (\cref{tab:mc}, Panel C), the test rejects
in 7.5\% of 200 null draws at a nominal 5\% level and in 10\%, 30\%, 38\%, 88\% and 100\% of draws as the mule share
rises from 0.25\% to 4\%. The mild over-rejection reflects a bootstrap standard error about 12\% smaller than the
Monte Carlo standard deviation; the plan therefore pairs the test with placebo thresholds. Because the standard error
scales with the square root of the number of deposits near the threshold, a platform with a hundred times the
simulated activity would detect a structuring share an order of magnitude smaller.

\begin{figure}[!htbp]
\centering
\fig{0.98}{fig_bunching}
\caption{Deposit amounts near the 50,000-baht face-verification threshold in one simulated exchange-year with 2\% mule
accounts, by the asset the deposit funds. Dashed lines are counterfactual densities fitted outside the shaded window;
coloured bars are the excess mass. \emph{Simulated data.}}
\label{fig:bunching}
\end{figure}

\paragraph{Speed cannot separate mules from honest fast traders}
\Cref{fig:dwellroc}(a) shows why. Arbitrageurs exit faster than most mules, and honest peer-to-peer merchants exit as
fast as camouflaged ones. A speed-only score reaches an AUC of 0.83 but a precision at $k$ of only 0.06--0.21
(\cref{tab:mc}), consistent with the bound of \cref{lem:speed}. Panel (b) shows the trade-off at the low false-positive
rates that matter in practice.

\begin{figure}[!htbp]
\centering
\fig{0.98}{fig_dwell_roc}
\caption{(a) Kernel densities of the time from deposit to on-chain withdrawal by population, logarithmic scale.
(b) ROC curves of five detectors in the same simulated exchange-year; the horizontal axis is logarithmic above 1\%.
\emph{Simulated data.}}
\label{fig:dwellroc}
\end{figure}

\paragraph{Direction matters}
The undirected two-hop rule flags every bridge trader and honest merchant, because their withdrawals go to exchange hot
wallets that sit, on average, 2.05 undirected hops from a flagged address (\cref{fig:galaxy}). Its precision at $k$ is
0.02--0.06 at every prevalence---the worst of all detectors---even though its AUC of 0.77 looks respectable. Following
funds forward removes this artefact: directed exposure reaches an AUC of 0.98. Its precision at $k$, however, rises
only from 0.02 to 0.86 as prevalence grows, because honest remitters who use one of the tainted merchants are fully
exposed and tie with mules at the top of the ranking. This is taint contagion, and it is the reason exposure should be
one input to a score rather than a verdict.

\begin{figure}[!htbp]
\centering
\fig{0.98}{fig_galaxy}
\caption{A sample of the simulated transaction graph. Nodes are placed on rings by their forward distance to the
nearest flagged address and grouped by type; blue lines are payments into exchange hot wallets. \emph{Simulated data.}}
\label{fig:galaxy}
\end{figure}

\paragraph{Layered evidence holds up}
Scores that combine behaviour and exposure dominate (\cref{fig:power} and \cref{tab:mc}). The isolation forest on all
indicators attains a precision at $k$ of 0.78 at the lowest prevalence and 0.90--0.92 from 1\% upwards; without the
exposure indicator it attains 0.70--0.83, so blockchain intelligence adds roughly eight to ten points of precision.
The transparent red-flag composite is less precise (about 0.60 at every prevalence) but each of its components can be
explained to an account holder or a court, a property that matters when accounts are frozen.

\begin{figure}[!htbp]
\centering
\fig{0.98}{fig_power}
\caption{Monte Carlo summary. (a) Precision at $k$ by detector and mule share (50 exchange-years per cell).
(b) Rejection rate of the bunching-difference test (200 draws under the null). (c) Rejection rate of synthetic control
by true effect size (200 country panels per effect). \emph{Simulated data.}}
\label{fig:power}
\end{figure}

\begin{table}[!htbp]
\centering\small
\begin{threeparttable}
\caption{Monte Carlo performance of detectors and of the bunching-difference test}
\label{tab:mc}
\begin{tabular}{lccccc}
\toprule
 & \multicolumn{5}{c}{Share of accounts that are mules} \\
\cmidrule(lr){2-6}
Detector & 0.25\% & 0.5\% & 1\% & 2\% & 4\% \\
\midrule
\multicolumn{6}{l}{\emph{Panel A. Area under the ROC curve}} \\
Undirected 2-hop exposure & 0.774 & 0.775 & 0.775 & 0.773 & 0.775 \\
Dwell time only & 0.826 & 0.827 & 0.828 & 0.829 & 0.827 \\
Directed 3-hop exposure & 0.981 & 0.981 & 0.980 & 0.981 & 0.982 \\
Red-flag composite, behaviour & 0.958 & 0.960 & 0.959 & 0.960 & 0.959 \\
Red-flag composite, all & 0.982 & 0.982 & 0.982 & 0.983 & 0.982 \\
Isolation forest, behaviour & 0.996 & 0.997 & 0.997 & 0.997 & 0.995 \\
Isolation forest, all & 0.999 & 0.999 & 0.999 & 0.999 & 0.999 \\
\addlinespace
\multicolumn{6}{l}{\emph{Panel B. Precision at $k$}} \\
Undirected 2-hop exposure & 0.02 & 0.02 & 0.03 & 0.05 & 0.06 \\
Dwell time only & 0.06 & 0.07 & 0.10 & 0.12 & 0.21 \\
Directed 3-hop exposure & 0.02 & 0.10 & 0.45 & 0.71 & 0.86 \\
Red-flag composite, behaviour & 0.58 & 0.60 & 0.60 & 0.60 & 0.60 \\
Red-flag composite, all & 0.59 & 0.61 & 0.60 & 0.60 & 0.60 \\
Isolation forest, behaviour & 0.70 & 0.76 & 0.80 & 0.83 & 0.81 \\
Isolation forest, all & 0.78 & 0.84 & 0.90 & 0.92 & 0.92 \\
\addlinespace
\multicolumn{6}{l}{\emph{Panel C. Bunching-difference test}} \\
Rejection rate at 5\% & 0.10 & 0.30 & 0.38 & 0.88 & 1.00 \\
Mean $\hat\beta_{\mathrm{USDT}}-\hat\beta_{\mathrm{other}}$ & 0.014 & 0.056 & 0.088 & 0.180 & 0.281 \\
\bottomrule
\end{tabular}
\begin{tablenotes}[flushleft]\footnotesize
\item \emph{Notes:} Means over 50 simulated exchange-years of 8,000 accounts per column. Precision at $k$ is the share of
mules among the $k$ highest-scored accounts, with $k$ equal to the number of mules. Under the null of no mules, the
bunching-difference test rejects in 7.5\% of 200 draws (mean difference $-0.007$). \emph{Simulated data.}
\end{tablenotes}
\end{threeparttable}
\end{table}

\paragraph{Synthetic control is well sized}
For policy shocks I simulate weekly flow indices for a treated country and 19 donors driven by two persistent common
factors, plant an effect after week 104, and apply synthetic control with rank-based placebo inference. In the example
of \cref{fig:synth}, the method recovers a planted effect of $-0.150$ log points as $-0.155$ with seven donors receiving
weight, and the treated unit ranks first among 20 ($p=0.05$). Across 200 panels per effect size (\cref{tab:scm}), the
false-rejection rate is 2.5\%, power is 38\% for a 0.05 log-point effect and 82\% for 0.10, and the mean estimate is
unbiased.

\begin{figure}[!htbp]
\centering
\fig{0.98}{fig_synth}
\caption{Synthetic control on a simulated country panel with a planted effect of $-0.15$ log points. (a) Treated unit
and its synthetic counterpart. (b) The treated gap (red) against the gaps of 19 placebo units. \emph{Simulated data.}}
\label{fig:synth}
\end{figure}

\begin{table}[!htbp]
\centering\small
\begin{threeparttable}
\caption{Size and power of synthetic control in simulated country panels}
\label{tab:scm}
\begin{tabular}{lcccccc}
\toprule
True effect (log points $\times$ 100) & 0 & 5 & 10 & 15 & 20 & 30 \\
\midrule
Rejection rate ($p \le 0.05$) & 0.025 & 0.380 & 0.820 & 0.970 & 0.995 & 1.000 \\
Mean estimated effect & -0.001 & -0.051 & -0.101 & -0.151 & -0.201 & -0.301 \\
\bottomrule
\end{tabular}
\begin{tablenotes}[flushleft]\footnotesize
\item \emph{Notes:} 200 panels per column, 20 units, 104 pre-treatment and 26 post-treatment weeks; rejection if the
treated unit has the largest post/pre RMSPE ratio ($p\le 0.05$). \emph{Simulated data.}
\end{tablenotes}
\end{threeparttable}
\end{table}

\subsection{What transfers to real data}

The absolute levels in \cref{tab:mc} depend on choices about populations and tradecraft that real data will revise.
Four conclusions do not depend on them. First, a statistic that compares the stablecoin channel with another channel on
the same platform is far more robust than one that measures the stablecoin channel alone, and its normalisation must
not depend on the shape of the amount distribution. Second, speed and undirected proximity, the two heuristics most
often described in compliance practice, fail by construction whenever honest fast traders and exchange hubs exist.
Third, at low prevalence AUC flatters detectors and precision at $k$ should be reported. Fourth, combining behaviour
with direction-aware blockchain intelligence is what makes detection precise; neither suffices alone.
